% Source: https://artofproblemsolving.com/community/c3h1486340_help_with_root_problem
% Stable identity: aops-topic-1486340-q01
Let $x=\sqrt{4+\sqrt7}-\sqrt{4-\sqrt7}$. Then $x>0$, because $4+\sqrt7>4-\sqrt7>0$. Squaring,
\[
x^2=(4+\sqrt7)+(4-\sqrt7)-2\sqrt{(4+\sqrt7)(4-\sqrt7)}=8-2\sqrt{16-7}=8-6=2 .
\]
So $x=\boxed{\sqrt2}\approx1.414$. Neither claim is right: $2$ is the value of $x^2$, not of $x$, and $2\sqrt7\approx5.29$ is far too large (the expression is less than $\sqrt{4+\sqrt7}<3$).

(Alternatively, $\sqrt{4\pm\sqrt7}=\sqrt{\dfrac{8\pm2\sqrt7}{2}}=\dfrac{\sqrt7\pm1}{\sqrt2}$, whose difference is $\dfrac{2}{\sqrt2}=\sqrt2$.)
